\documentclass[10pt,letterpaper]{article}
\usepackage{cvstyle}
\usepackage[french]{babel}

\begin{document}

\cvheader
  {Marc-Antoine Dupuis}
  {Étudiant à la maîtrise en science politique | Université Laval}
  {\href{mailto:MADUP212@ulaval.ca}{MADUP212@ulaval.ca}\cvsep
   \href{https://www.linkedin.com/in/marc-antoine-dupuis-a6a069233/}{LinkedIn}\cvsep
   \href{https://www.masdupuis.com/}{masdupuis.com}}

\section*{Formations}
\cventry{Depuis 2025}
  {\textbf{Université Laval, Québec | Directeur : Yannick Dufresne, Ph.\,D.}, \textit{Maîtrise en science politique}}{}
\cventry{2021--2025}
  {\textbf{Université d'Ottawa (2021--2023) et UQÀM (2023--2025)}, \textit{Baccalauréat en science politique}}{}

\section*{Recherche en cours}
\cventry{Depuis 2026}
  {\textbf{Projet de maîtrise, Université Laval}, \textit{Maîtrise par article}}{}
\begin{cvitems}
\item Dupuis, M. et Dufresne, Y. (s.\,d.). \textit{Comment les classes sociales se sont-elles reconfigurées
      au Québec selon les groupes linguistiques ?} [soumis à SPSA 2027].
\item Dupuis, M. et Dufresne, Y. (s.\,d.). \textit{Variables sociodémographiques et sélection des candidatures
      aux élections générales québécoises de 2026 : une analyse tenant compte de la compétitivité électorale.}
\end{cvitems}

\section*{Projets en cours}
\cventry{Depuis 2026}
  {\textbf{Prof. Datagotchi vous répond}, \textit{Centre d'analyse des politiques publiques}}
  {Co-fondateur. Assistant conversationnel (RAG) qui répond, dans la peau du Prof. Datagotchi, aux questions
   sur l'élection générale québécoise de 2026, à partir des entrevues, des plateformes et des communiqués des
   partis. \href{https://prof-datagotchi.com/}{prof-datagotchi.com}}
\cventry{Depuis 2026}
  {\textbf{Polimètre Mark Carney}, \textit{Centre d'analyse des politiques publiques}}
  {Suivi des promesses électorales du gouvernement fédéral.}

\section*{Projets passés}
\cventry{2026}
  {\textbf{Boussole électorale québécoise}, \textit{Vox Pop Lab}}{}
\cventry{2026}
  {\textbf{Boussole électorale, élections municipales en Ontario}, \textit{Vox Pop Lab}}{}
\cventry{2026}
  {\textbf{Polimètre Legault-Fréchette}, \textit{Centre d'analyse des politiques publiques}}{}
\cventry{2026}
  {\textbf{Polimètre Trudeau}, \textit{Centre d'analyse des politiques publiques}}{}
\cventry{2025}
  {\textbf{Projet de recherche sur le 2\textsuperscript{e} bataillon de construction}, \textit{UQÀM}}
  {Recherche portant sur la Première Guerre mondiale.}

\section*{Expériences professionnelles}
\cventry{Depuis 2021}
  {\textbf{Forces armées canadiennes (réserve)}, \textit{Administrateur des ressources humaines}}{}
\cventry{Depuis 2026}
  {\textbf{Université Laval}, \textit{Auxiliaire d'enseignement}}
  {Automne 2026 :}
\begin{cvitems}
\item École interdisciplinaire outils et méthodes (EIOM)
\item POL-2425, Persuasion et opinion publique
\item POL-6078, Outils numériques en sciences sociales
\item POL-2000, Méthodologie quantitative
\end{cvitems}
\cventry{2024--2025}
  {\textbf{MIFI}, \textit{Auxiliaire de recherche et aide à la coordination}}{}

\section*{Conférences et enseignement}
\cventry{2026}
  {\textbf{Introduction à la revue de littérature}, \textit{POL-6078, Université Laval}}{}
\cventry{2026}
  {\textbf{Atelier sur les langages de balisage}, \textit{EIOM, Université Laval}}{}
\cventry{2026}
  {\textbf{Atelier sur les fondamentaux de LaTeX}, \textit{Université Laval}}{}
\cventry{2026}
  {\textbf{Modérateur de panel}, \textit{UNIDEF 2026}}{}

\section*{Groupes de recherche}
\cventry{Depuis 2026}
  {\textbf{Chaire de leadership en enseignement numérique (CLESSN), Université Laval}, \textit{Membre}}
  {Chercheur étudiant.}
\cventry{Depuis 2026}
  {\textbf{Centre d'analyse des politiques publiques (CAPP), Université Laval}, \textit{Membre}}
  {Chercheur étudiant.}
\cventry{Depuis 2026}
  {\textbf{Groupe de recherche en communication politique (GRCP), Université Laval}, \textit{Membre}}
  {Chercheur étudiant.}
\cventry{Depuis 2026}
  {\textbf{Centre pour l'étude de la citoyenneté démocratique (CÉCD), Université Laval}, \textit{Membre}}
  {Chercheur étudiant.}

\section*{Bourses et reconnaissances}
\cventry{2026}{\textbf{Centre sur la sécurité internationale}, \textit{Bourse de mérite}}{}
\cventry{2026}{\textbf{Centre d'analyse des politiques publiques}, \textit{Bourse de mérite}}{}
\cventry{2024}{\textbf{UQÀM}, \textit{Bourse de mobilité}}{}
\cventry{2022}{\textbf{Université d'Ottawa}, \textit{Bourse de mérite}}{}

\section*{Compétences spécifiques}
\cvskill{Langages de programmation}{R, HTML, CSS, LaTeX, Markdown}
\cvskill{Outils et logiciels}{Git/GitHub, Notion, Microsoft Office}
\cvskill{Langues}{Français (langue maternelle) et anglais (courant)}

\section*{Références}
Fournies sur demande.

\end{document}
